% !TeX root = ../../Book3.tex
\section{函数方程}
\label{b3:sec:2604}

延拓积分的两个幂次在 \(s\) 与 \(1-s\) 互换时交换位置。这是函数方程的来源，而不是由零点图形推测出的对称。

\subsection{\texorpdfstring{\(\theta\)}{\AnalysisPdfTheta} 变换}
\label{b3:subsec:260401}

在\eqref{b3:eq:26-completed-integral}中，显式有理项和积分都在 \(s\mapsto1-s\) 下不变。所有换元都最初在绝对收敛域完成，而最终的等式是全平面的亚纯恒等式；不存在要求 \(s\) 与 \(1-s\) 同时落在原 Dirichlet 收敛域的矛盾。

\subsection{完备 \texorpdfstring{\(\zeta\)}{\AnalysisPdfZeta} 函数}
\label{b3:subsec:260402}

定义\term[wanbei-zeta-hanshu]{完备 zeta 函数}[completed zeta function]
\(\Lambda_\zeta(s)\coloneqq\pi^{-s/2}\Gamma(s/2)\zeta(s)\)。下标用来区别算术函数 \(\Lambda(n)\)。它只在零与一有单极点，留数分别为负一与一，并满足
\begin{equation}\label{b3:eq:26-functional-equation}
\Lambda_\zeta(s)=\Lambda_\zeta(1-s).
\end{equation}
\symidx{\(\Lambda_\zeta(s)\)：完备 zeta 函数}

\subsection{函数方程}
\label{b3:subsec:260403}

利用 Gamma 的反射和倍增公式，将\eqref{b3:eq:26-functional-equation}中的 Gamma 比值化简，得到另一形式
\[
\zeta(s)=2^s\pi^{s-1}\sin(\pi s/2)\Gamma(1-s)\zeta(1-s).
\]
化简可从 \(\Gamma(s/2)\Gamma(1-s/2)=\pi/\sin(\pi s/2)\) 和
\(\Gamma((1-s)/2)\Gamma(1-s/2)=2^s\sqrt\pi\Gamma(1-s)\) 出发。式中各因子在某些整数处分别无穷或为零，等式应按亚纯延拓解释，不把那些点的形式乘积直接当作普通数值计算。

\subsection{对称性}
\label{b3:subsec:260404}

在讨论零点以前，先以两个局部计算检验这些因子的配合。令 \(s\to0\)，由已知 Laurent 展开和 Gamma 在一处的导数，
\[
 \zeta(1-s)=-\frac1s+\gamma+O(s),\qquad
 \Gamma(1-s)=1+\gamma s+O(s^2).
\]
两者乘积的常数项恰好抵消。再展开
\[
 2^s\pi^{s-1}=\pi^{-1}\bigl(1+s\log(2\pi)+O(s^2)\bigr),
 \qquad \sin(\pi s/2)=\frac{\pi s}2+O(s^3),
\]
得到
\[
 \zeta(s)=-\frac12-\frac12\log(2\pi)s+O(s^2).
\]
所以 \(\zeta'(0)=-\tfrac12\log(2\pi)\)。这里有限结果来自正弦零点与 \(\zeta(1-s)\) 极点的抵消，不是把两个无定义的点值直接相乘。

\ProseParagraphBreak

另一方面，对整数 \(m\ge1\)，在 \(s=1-2m\) 处各因子都是普通有限值，因而
\[
 \zeta(1-2m)=(-1)^m\frac{2(2m-1)!}{(2\pi)^{2m}}\zeta(2m).
\]
例如 \(\zeta(-1)=-1/12\)。负奇数值由正偶数值确定，而负偶数处由正弦因子产生零点，两种情形不可混淆。以上数值是延拓函数的值，不是把原发散 Dirichlet 级数作通常求和得到的值。

\ProseParagraphBreak

因为原级数实系数，\(\zeta(\overline s)=\overline{\zeta(s)}\)，延拓保持此式。完备函数还有 \(s\mapsto1-s\) 对称，所以非实非平凡零点按 \(\rho,\overline\rho,1-\rho,1-\overline\rho\) 出现；四点可能重合为两点，例如落在中线时。这个对称关系限制零点位置，却不强迫每条对称轨道都缩在中线上。
